% !TEX root = main.tex
% !TEX root = main.tex
\begin{table*}[!th]
\resizebox{\textwidth}{!}{
\begin{threeparttable}
\caption{Comparison of PCS Performance on $8$-bit $2^{30}$-Vectors}
\label{table:experimentPCS}
\begin{tabular}
	{ >{\arraybackslash}p{2.1cm}
	>{\centering\arraybackslash}p{1.6cm} 
	>{\centering\arraybackslash}p{1.6cm} 
	>{\centering\arraybackslash}p{1.6cm}
	>{\centering\arraybackslash}p{1.8cm}
	>
    %add a line to separate those run out of memory
    {\centering\arraybackslash}p{1.5cm}
	>{\centering\arraybackslash}p{1.5cm}
 >{\centering\arraybackslash}p{2.0cm} 
	>{\centering\arraybackslash}p{1.5cm} 
	}
\hline
	&	\makecell{KZG10\cite{asiacrypt/KateZG10}} 
	& \makecell{Sonic\cite{ccs/MallerBKM19}}
	& \makecell{Marlin\cite{eurocrypt/ChiesaHMVW20}} 
	& \makecell{BCMS20\cite{tcc/BunzCMS20}}
	& \makecell{Hyrax\cite{sp/WahbyTSTW18}}
	& \makecell{Ligero\cite{ccs/AmesHIV17}}
	% \makecell{Libra on \\ $512 \times 512$ matrices\tnote{c} } 
	& \makecell{Brakedown\cite{crypto/GolovnevLSTW21}}
	% \makecell{Spartan with \\ $2^{24}$ constraints\tnote{d} }
	& \makecell{This Work}
	% \makecell{HyperPlonk\tnote{d} with \\ $2^{24}$ constraints\tnote{d} } 
	% & \makecell{Ratio}
\\
\hline
Setup time 
& $933.38$s 
& $3,837.30$s
& $814.80$s 
& $1,428.84$s
& $13.35$ms
& $1.36\mu$s 
& $2.59$s
& $1.02$s
\\
\hline
Public pars. size 
& $103.08$GB 
& $206.16$GB
& $206.16$GB 
& $137.43$GB 
& $2.10$GB
& $17$B 
& $688.85$MB
& $11.01$GB
\\
\hline
Commit time 
& $686.34$s
& $742.78$s
& $822.22$s  
& $305.94$s
& $171.03$s
& $130.32$s 
& $127.03$s
& $192.90$s
\\
\hline
Commit size 
& $48$B 
& $96$B
& $97$B 
& $65$B
& $2.10$
& $64$B
& $64$B
& $576$B
\\
\hline
Open time 
& $967.14$s  
& $1,019.46$s
& $1,020.11$s  
& $4,647.62$s
& $3.72$s %
& $98.27$s % 
& $74.21$s
& $5.77$s
\\
\hline
Proof size 
& $81$B
& $57$B
& $57$B 
& $2.07$KB
& $1.05$MB % 
& $57.78$MB 
& $215.59$MB
& $36.78$KB
\\ 
\hline
Verifier time 
& $2.98$ms 
& $1.94$ms 
& $3.13$ms 
& $1,897.38$s 
& $1.21$s %
& $4.92$s
& $22.37$s
& $56.10$ms
\\
\hline
Peak RAM 
& $789.08$GB  
& $1,034.21$GB
& $818.22$GB 
& $1,154.64$GB
& $102.12$GB %
& $202.61$GB % 
& $187.89$GB
& $3.46$GB
\\ 
\hline
\end{tabular}
\begin{tablenotes}[para,flushleft]
    Parallelized over 64 threads of CPUs.
    Due to memory overflow on $2^{30}$-vectors, the performance of KZG10, Sonic, Marlin, and BCMS20 was estimated based on results for smaller $2^{29}$- or $2^{28}$-vectors with corresponding scaling factors.
\end{tablenotes}
\end{threeparttable}
}
\end{table*}

\section{Base PCS Optimization}
\label{appx:pcs}

In our experiments, we employ a variant of Dory~\cite{tcc/Lee21} as the base PCS.
We compare the performance of this PCS with other PCSs implemented in the
\textsf{arkworks} library, and report the results in
Table~\ref{table:experimentPCS}, for committing to and opening a
size-$2^{30}$ 8-bit vector.

The total commitment-plus-opening time is comparable to the measured Brakedown
implementation, while our PCS has smaller proof size and verifier time on this
benchmark and retains logarithmic asymptotic proof size and verifier time.

The PCS does not convert an entire 8-bit vector into field elements during
commitment or opening.  Its $3.46$\,GB peak RAM is therefore more than an order
of magnitude below every other implementation in
Table~\ref{table:experimentPCS}.

Our PCS is compatible with bucket-wise computation.  For a BVA, both commitment
and opening iterate over the original vector entries and their logical offsets;
neither intra-bucket zero suffixes nor the final implicit zero suffix is materialized.

Below, we briefly describe the construction of this PCS.

\subsection{Pairing Related Notations}

\mysmallskip {\bf Pairing.}
\label{sec:notationPairing}
A pairing is a bilinear map $e: \Group_1 \times \Group_2 \mapsto \Group_T$ over three cyclic groups $\Group_1$, $\Group_2$, and $\Group_T$, each of prime order $p$, such that:
\[
	e(aA, bB) = ab\, e(A, B), \, \forall\, a, b\in \Field, \, A \in \Group_1, B \in \Group_2,
\]
where $\Field$ denotes the ring of integers modulo $p$, and $\Field^* := \Field \backslash \{ 0 \}$.
We use the asymmetric (type-3) pairing, characterized by no efficiently computable homomorphism between $\Group_1$ and $\Group_2$.

\mysmallskip{\bf Group elements.}
\label{sec:notationGroupElements}
Group elements such as $G, H, U, L, R, A, B$ are capitalized.
Commitments to $a, b, c, d \in \Field$ are denoted by 
$C_a, C_b, C_c, C_d \in \Group_T$, respectively.
We adopt the additive notation for group operations, using $\gadd$ for group addition and $\gzero$ for the identity element in $\Group_T$.
The summation of indexed group elements is represented as $\gsum_{i = \mathit{start}}^{\mathit{end}}$.

\mysmallskip{\bf Group matrices and vectors.}
Group matrices or vectors are denoted by uppercase bold letters, \eg, $\bm{A}, \bm{B}, \vec{\bm{A}}, \vec{\bm{B}}$.

\mysmallskip{\bf Group matrix multiplication.}
Define the product of group elements $A \in \Group_1$ and $B \in \Group_2$ as their pairing $e(A, B)$:
\[
	AB := BA := e(A, B).
\]
This allows for extending matrix multiplication to group matrices in a manner analogous to field matrices.
Consider $\bm{A} = \{ A_{ik}\} \in \Group_1^{m\times l}$ and $\bm{B} = \{ B_{kj}\} \in \Group_2^{l \times n}$, their product $\bm{A} \bm{B}$ is an $m\times n$ matrix with the $(i, j)$ entry defined as 
$\gsum_{k = 1}^l e(A_{ik}, B_{kj})
$.
Multiplying a group element $A \in \Group_1$ by a group matrix $\bm{B} = \{ B_{kj}\} \in \Group_1^{l \times n}$ results in a group matrix $\{ e(A, B_{kj}) \} $.
When $\bm{B} = \{ B_{ik}\} \in \Group_2^{m \times l}$ and $\bm{A} = \{ A_{kj}\} \in \Group_1^{l \times n}$, their product $\bm{B}\bm{A}$ is also well defined.
Matrix multiplication between field matrices and group matrices is defined similarly.


\mysmallskip{\bf Inner product.}
The inner product is a special case of matrix multiplication.
The inner product of two group vectors $\vec{\bm{A}} \in \Group_1^n$ and $\vec{\bm{B}} \in \Group_2^n$ is defined as:
\[
	\langle \vec{\bm{A}}, \vec{\bm{B}} \rangle :=
	\vec{\bm{A}}^\top \vec{\bm{B}} = 
	\vec{\bm{B}}^\top \vec{\bm{A}} = 
	\langle \vec{\bm{B}}, \vec{\bm{A}} \rangle .
\]


\mysmallskip {\bf Partition.}
Given a vector $\vec{\bm{A}} \in \Group_T^n$, we define its left half as 
$\vec{\bm{A}}_L := (A_1, \ldots, A_{n/2})$ and its right half as 
$\vec{\bm{A}}_R := (A_{n/2 + 1}, \ldots, A_{n})$.
This partitioning procedure is denoted by
$(\vec{\bm{A}}_L || \vec{\bm{A}}_R) \Half \vec{\bm{A}}$.
For a matrix $\bm{a} \in \Group_T^{m\times n}$, it can be divided into left and right halves denoted as $\bm{A}_L$ and $\bm{A}_R$,
or into upper and lower halves represented by $\bm{A}_U$ and $\bm{A}_D$.


\subsection{Cryptographic Assumptions}
\label{appx:assumptions}
%Our definitions for Pedersen commitment, Pedersen vector commitment, and discrete logarithm relation are the same as the literature~\cite{sp/BuenzBBPWM18}.

Our $\Group_T$-Bulletproofs rely on the following $q$-type assumptions:
\begin{assumption}[$q$-PDH~\cite{asiacrypt/Groth10,cacm/ParnoHGR16}]
\label{def:pdh}
	The $q$-power Diffie-Hellman ($q$-PDH) assumption holds for a group $\Group$ if for all PPT adversary $\Adversary$:
	\begin{align}
		%\label{eq:qPKE}
		& \Pr
		\left(
			\left.
			\begin{aligned}
				& C = \hat{s}^{q + 1}\hat{G} \\ 
			\end{aligned}
			\right.
			\left|
			\begin{aligned}
				& (p, \Group) \leftarrow \Setup(1^{\lambda}); \\
				& \hat{G} \ChooseRandom \Group; \hat{s}\ChooseRandom \Field^*; \\
				& \srs \leftarrow (p, \Group; [\hat{G}], [\hat{s}\hat{G}], \ldots, [\hat{s}^q \hat{G}]);\\
				& C \leftarrow \Adversary(\srs)
			\end{aligned}
			\right.
		\right)
		\le \Negl(\lambda).
		\nonumber
	\end{align}
\end{assumption}


\begin{assumption}[$q$-PKE~\cite{asiacrypt/Groth10,cacm/ParnoHGR16}]
	The $q$-power knowledge of exponent ($q$-PKE) assumption holds for $(\Group_1, \Group_2, \Group_T, e)$ if for all PPT adversary $\Adversary$, there exists a non-uniform PPT extractor $\Extractor_{\Adversary}$ such that:
	\begin{align}
		% \label{eq:qPKE}
		& \Pr
		\left(
			\left.
			\begin{aligned}
				& V' = {\nu}V \\
				& \wedge V \ne (\sum_{i = 0}^{q} {\alpha}_i \hat{s}^i)\hat{G}\\ 
			\end{aligned}
			\right.
			\left|
			\begin{aligned}
				& (p, \Group_1, \Group_2, \Group_T, e) \leftarrow \Setup(1^{\lambda}); \\
				& \hat{G} \ChooseRandom \Group_i, i\in \{1, 2\}; \nu, \hat{s}\ChooseRandom \Field^*; \\
				& \srs \leftarrow (p, \Group_1, \Group_2, \Group_T, e; \hat{G}, \hat{s}\hat{G}, \ldots, \\
				& \quad \hat{s}^q \hat{G}, \nu \hat{G}, \nu \hat{s}\hat{G}, \ldots, \nu \hat{s}^q \hat{G});\\
				& (V', V; \alpha_0, \ldots, \alpha_q) \leftarrow (\Adversary|| \Extractor_{\Adversary})(\srs, \bm{z})
			\end{aligned}
			\right.
		\right)
		\nonumber \\ 
		& \le \Negl(\lambda),
		\nonumber
	\end{align}
for any auxiliary information $\bm{z} \in \{0, 1\}^{\Poly(\lambda)} $ that is generated independently of $\nu$.
Note that $(y_1; y_2) \gets (\Adversary|| \Extractor_{\Adversary})(x)$ means that on input $x$, $\Adversary$ outputs $y_1$ and
that $\Extractor_{\Adversary}$ with the same input $x$ and $\Adversary$'s random tape, produces $y_2$.
%and for the class Z of auxiliary input generators for every non-uniform PPT auxiliary input generator $Z\in Z$ and
\end{assumption}

%~\cite{asiacrypt/KateZG10}
\begin{assumption}[$q$-BSDH]
	The $q$-bilinear strong Diffie-Hellman ($q$-BSDH) assumption holds for $(\Group_1, \Group_2, \Group_T)$ if for all PPT adversary~$\Adversary$:
	\begin{align}
		% \label{eq:qPKE}
		& \Pr
		\left(
			\left.
			\begin{aligned}
				& E = (\hat{s} - r)^{-1} e(\hat{G}, \hat{H}) \\ 
			\end{aligned}
			\right.
			\left|
			\begin{aligned}
				& (p, \Group, \Group_T, e) \leftarrow \Setup(1^{\lambda}); \\
				& \hat{G} \ChooseRandom \Group_1; \hat{H} \ChooseRandom \Group_2; \hat{s}\ChooseRandom \Field^*; \\
				& \srs \leftarrow (p,\Group, \Group_T, e; \hat{G}, \hat{s}\hat{G}, \ldots, \\
				& \quad \hat{s}^q \hat{G}, \hat{H}, \hat{s}\hat{H}, \ldots, \hat{s}^q \hat{H})\\
				& (E \in \Group_T, r \in \Field^*) \leftarrow \Adversary(\srs)
			\end{aligned}
			\right.
		\right)
		\nonumber \\
		& \le \Negl(\lambda) .
		\nonumber
	\end{align}
\end{assumption}

\subsection{Two-Tier Commitment}
\label{appx:twotier}
The two-tier commitment~\cite{asiacrypt/Groth11} to a matrix $\bm{a} \in \Field^{n \times n}$ is defined as:
\begin{itemize}
\item \textbf{First-tier commitment:}
Apply the Pedersen vector commitment to each column or row of $\bm{a}$ using the base vector $\vec{\bm{G}} \in \Group_1^n$ or $\vec{\bm{H}} \in \Group_2^n$.
This results in $\vec{\bm{A}}_G \in \Group_1^n$ or $\vec{\bm{A}}_H \in \Group_2^n$ as follows:
	\[
		\vec{\bm{A}}_G = \vec{\bm{G}}^\top\bm{a} \quad 
		\text{or} \quad
		\vec{\bm{A}}_H = \bm{a}\vec{\bm{H}}.
	\]
\item \textbf{Second-tier commitment:}
Commit to $\vec{\bm{A}}_G$ or $\vec{\bm{A}}_H$ using AFGHO commitments~\cite{crypto/AbeFGHO10}, resulting in $C_a \in \Group_T$:
	\[
		C_a = \langle \vec{\bm{A}}_G, \vec{\bm{H}} \rangle = \langle \vec{\bm{A}}_H, \vec{\bm{G}} \rangle 
		= \vec{\bm{G}}^\top\bm{a}\vec{\bm{H}}.
	\]
We can add a hiding factor $\tilde{a} \in \Field$ with the base $\tilde{U} \in \Group_T$ to the two-tier commitment as follows:
	\[
		C_a = \vec{\bm{G}}^\top\bm{a}\vec{\bm{H}} \, \gadd \, \tilde{a} \tilde{U}.
	\]
\end{itemize}

\mysmallskip
{\bf Caching first-tier commitment and other optimizations.}
In the commitment phase, the prover caches the first-tier commitments $\vec{A}_G \in \Group_1^n$ or $\vec{A}_H \in \Group_2^n$.
These caches are reused in the PCS opening.

\mysmallskip
{\bf Bucket-wise commitment.}
For the expanded vector produced by a BVA, the two-tier commitment can be
computed bucket by bucket without materializing padded coordinates.  The
expanded vector is conceptually reshaped into a square matrix and committed
using the two-tier commitment scheme.  During the first-tier commitment, the
prover iterates only over stored input entries, computes the logical position of
each entry, multiplies it by the corresponding group element in $\vec{\bm{G}}$,
and accumulates the result into the first-tier commitment.

\subsection{Bulletproofs in pairing group.}
\mysmallskip
{\bf Trusted Setup.}
We set up the below elements in $(\Group_1, \Group_2, \Group_T, e)$,
with the hyperparameter $q$ exceeding the largest supported matrix dimension.
$\hat{s}$ and $\nu$ 
%are the ``toxic waste'' that 
must be securely discarded after setup.
\(
\begin{aligned}
	&
	\hat{s}, \nu \ChooseRandom \Field^*; 
	\quad
	\hat{G} \ChooseRandom \Group_1;
	\,\,
	\hat{H} \ChooseRandom \Group_2;
	\,\,
	G' \leftarrow \hat{s}\hat{G} \in \Group_1;
	\,\,
	H' \leftarrow \hat{s}^q\hat{H} \in \Group_2;
	\\
 	& \hat{G}'' \leftarrow \nu \hat{G}; \,\, \hat{H}'' \leftarrow \nu \hat{H}; \quad
 	\vec{\bm{G}} \leftarrow (\hat{G}, \hat{s}\hat{G}, \ldots, \hat{s}^{q - 1}\hat{G});
	\\
 	&
 	\vec{\bm{H}} \leftarrow (\hat{H}, \hat{s}^q\hat{H}, \ldots, \hat{s}^{(q - 1)q}\hat{H}); \quad
	 \vec{\bm{G}}' \leftarrow (\hat{G}, \hat{s}^q\hat{G}, \ldots, \hat{s}^{(q - 1)q}\hat{G});
	\\
 	&
 	\vec{\bm{G}}'' \leftarrow (\nu \hat{G}, \nu \hat{s}\hat{G}, \ldots, \nu \hat{s}^{q - 1}\hat{G}); \,
 	\vec{\bm{H}}'' \leftarrow (\nu \hat{H}, \nu \hat{s}^q\hat{H}, \ldots, \nu \hat{s}^{(q - 1)q}\hat{H});
	\\
 	&
 	G_0 \leftarrow \hat{s}^{-1} \hat{G}; \quad H_0 \leftarrow \hat{H};\quad 
 	U \leftarrow \hat{s}^{-1} \hat{G} \hat{H}; \quad \tilde{U} \leftarrow \hat{s}^{-2} \hat{G} \hat{H}.
\end{aligned}
\)

This setup is similar to that in BMMTV21~\cite{asiacrypt/BuenzMMTV21}.
With this setup, an inner product $C = \vec{\bm{\xi}} \bullet \vec{\bm{G}}$ for public $\vec{\bm{\xi}} \in \Zp$ is equivalent to $C = \phi(\hat{s})\hat{G}$, where $\phi$ is an $n$-degree univariate polynomial with coefficients $\vec{\bm{\xi}}$.
This polynomial structure enables efficient pairing-based proofs for the inner product $C$.

In our trusted setup,
\(
 	\gsum_{i=1}^q G'_i H_{q+1-i} = \gzero
\)
forms a non-trivial inner-pairing product, 
violating a key assumption of Dory~\cite[Lemma~1]{tcc/Lee21}
for the absence of non-trivial inner-pairing-product relations.
To overcome this, we derive our PCS from Bulletproofs over $\Group_T$.
From the $(q^2 + 2)$-PDH assumption (Appendix~\ref{appx:assumptions}),
the assumption of no non-trivial discrete logarithm relation~\cite[Definition~1]{sp/BuenzBBPWM18} 
%for Bulletproofs
still holds 
%for our trusted setup 
among the $q^2$ entries of
\(
 	\vec{\bm{G}} \kronecker \vec{\bm{H}} \in \Group_T^{q^2},
\)
where the tensor product of group vectors is defined similarly to that of field vectors,
using pairings to compute group element products.

\iffalse
\mysmallskip
{\bf Two-tier commitment in $\Group_T$ as Pedersen vector commitment.}
Using the $\Group_T$ base $\vec{\bm{G}} \kronecker \vec{\bm{H}}$,
the Pedersen vector commitment~\cite[Definition~7]{sp/BuenzBBPWM18} for a flattened matrix $\bm{a} \in \Zp^{n \times n}$
aligns with the two-tier commitment~\cite{asiacrypt/Groth11},
which first commits to each column of $\bm{a}$ using $\vec{\bm{G}}$ to derive $\vec{\bm{A}} \in \Group_1$ and then apply the AFGHO commitments~\cite{crypto/AbeFGHO10}
to compute $C_a = \vec{\bm{A}} \bullet \vec{\bm{H}}$.
Specifically:
\[
 	\vectorize{\bm{a}} \bullet (\vec{\bm{G}} \kronecker \vec{\bm{H}}) = \vec{\bm{G}}^\top \bm{a} \vec{\bm{H}} = \vec{\bm{A}} \bullet \vec{\bm{H}} := C_a \in \Group_T.
\]
\fi


% !TEX root = main.tex
\begin{algorithm}[!t]
	\caption{Protocol of $\Group_T$-Bulletproofs}
	\label{alg:pcs}
	\KwIn{ Public input: $C \in \Group_T, \vec{\bm{l}} \in \Field^n, \vec{\bm{r}} \in \Field^n $ \;\\
					\qquad 
					$ \vec{\bm{G}}\in \Group_1^{n}, \vec{\bm{H}} \in \Group_2^{n}$; $G_0 \in \Group_1, H_0 \in \Group_2$;
	\\
	$ \qquad \Prover$'s private input: $\bm{a} \in \Field^{n \times n}$ and cache: \\
		$\qquad \qquad \qquad \vec{\bm{A}} = \vec{\bm{G}}^\top \bm{a} \in \Group_1^n$ }
% It is not closed because it is a G1 element.
%Left open if it is in G2.
%Closed if it is in GT.
	\KwOut{$\TRUE$ or $\FALSE$}
	$\Prover$ computes: 
	$\vec{\bm{v}} \leftarrow \vec{\bm{l}}^\top\bm{a} \in \Field^n$\;
    \label{line:computeal}
	\ForEach{$j \in [1, \log_2 n]$ }{
		$\Prover$ computes:\\
		\Indp
		$(\vec{\bm{A}}_\mathsf{L} || \vec{\bm{A}}_\mathsf{R}) \Half \vec{\bm{A}}$;\,
		$(\vec{\bm{H}}_\mathsf{L} || \vec{\bm{H}}_\mathsf{R}) \Half \vec{\bm{H}}$\;
		$(\vec{\bm{v}}_{\mathsf{L}} || \vec{\bm{v}}_{\mathsf{R}}) \Half \vec{\bm{v}}$;\,
		$(\vec{\bm{r}}_\mathsf{L} || \vec{\bm{r}}_\mathsf{R}) \Half \vec{\bm{r}}$\;
		$ L_{(j)} \leftarrow \langle \vec{\bm{A}}_{L}, \vec{\bm{H}}_\mathsf{R}\rangle \gadd
		\langle \vec{\bm{v}}_{L}, \vec{\bm{r}}_\mathsf{R} \rangle \cdot G_0H_o	
		$ \;
	\label{line:computeL1}
		$ R_{(j)} \leftarrow \langle \vec{\bm{A}}_{R}, \vec{\bm{H}}_\mathsf{L}\rangle \gadd
		\langle \vec{\bm{v}}_{R}, \vec{\bm{r}}_\mathsf{L}\rangle \cdot G_0H_0 
		$ \;
		\Indm
		$\PtoV$: $L_{(j)}, R_{(j)} \in \Group_T$\;
		$\VtoP$: $ x_{(j)} \ChooseRandom \Field^*$ \;
		$\Prover$ computes: \\
		\Indp
		$\vec{\bm{A}} \leftarrow \vec{\bm{A}}_\mathsf{L} \gadd x_{(j)}^{-1} \vec{\bm{A}}_\mathsf{R} $;
		$\vec{\bm{H}} \leftarrow \vec{\bm{H}}_\mathsf{L} \gadd x_{(j)} \vec{\bm{H}}_\mathsf{R} $\;
	\label{line:computeAH}
		$\vec{\bm{v}} \leftarrow \vec{\bm{v}}_{\mathsf{L}} + x_{(j)}^{-1} \vec{\bm{v}}_{\mathsf{R}} $;
		$\vec{\bm{r}} \leftarrow \vec{\bm{r}}_\mathsf{L} + x_{(j)} \vec{\bm{r}}_\mathsf{R} $\;
	\label{line:computevr}
		\Indm
	}
	% \tcp{Lengths of $\vec{\bm{A}}, \vec{\bm{v}}$, and $\vec{\bm{H}}$ become 1} 
	$\Prover$ and $\Verifier$ compute:\\
		\Indp
		$\vec{\bm{\xi}}_H \leftarrow (1, x_{(1)}) \kronecker \cdots \kronecker (1, x_{(\log_2 n)}) \in \Field^n$\;
	\Indm
	$\Prover$ computes:\\
		\Indp
	$\vec{\bm{a}} \leftarrow \bm{a} \vec{\bm{\xi}}_H^{-1} \in \Field^n$;\,	$V_H \leftarrow \vec{\bm{H}} \in \Group_1$;\, $r\leftarrow \vec{\bm{r}} \in \Field$\;
    \label{line:computear}
	\Indm
	\ForEach{$j \in [\log_2 n +1, 2\log_2 n]$ }{
	$\Prover$ computes:\\
	\Indp
		$(\vec{\bm{a}}_\mathsf{L} || \vec{\bm{a}}_\mathsf{R}) \Half \vec{\bm{a}}$;
		$(\vec{\bm{G}}_\mathsf{L} || \vec{\bm{G}}_\mathsf{R}) \Half \vec{\bm{H}}$;
		$(\vec{\bm{l}}_\mathsf{L} || \vec{\bm{l}}_\mathsf{R}) \Half \vec{\bm{l}}$\;
		$ L_{(j)} \leftarrow \langle \vec{\bm{a}}_{L}, \vec{\bm{G}}_\mathsf{R}\rangle \cdot V_H \gadd
        \langle \vec{\bm{a}}_{L}, \vec{\bm{l}}_\mathsf{R}\rangle
        \cdot (rG_0H_0)
        $\;
	\label{line:computeL2}
		$ R_{(j)} \leftarrow \langle \vec{\bm{a}}_{R}, \vec{\bm{G}}_\mathsf{L}\rangle \cdot V_H \gadd
		\langle \vec{\bm{a}}_{R}, \vec{\bm{l}}_\mathsf{L}\rangle
        \cdot (rG_0H_0)
		$\;
	\Indm
		$\PtoV$: $ L_{(j)}, R_{(j)} \in \Group_T$\;
		$\VtoP$: $ x_{(j)} \ChooseRandom \Field^*$ \;
	$\Prover$ computes:\\
	\Indp
		$\vec{\bm{a}} \leftarrow \vec{\bm{a}}_\mathsf{L} + x_{(j)}^{-1} \vec{\bm{a}}_\mathsf{R}, \vec{\bm{l}} \leftarrow \vec{\bm{l}}_\mathsf{L} + x_{(j)} \vec{\bm{l}}_\mathsf{R} \in \Field^{n/2^j}$\;
		$\vec{\bm{G}} \leftarrow \vec{\bm{G}}_\mathsf{L} \gadd x_{(j)} \vec{\bm{G}}_\mathsf{R} \in \Group_1^{n/2^j}$\;
	\label{line:computeG}
	\Indm
	}
	% \tcp{Now, the length of $\vec{\bm{a}}$ is 1}
		$\PtoV$: $\hat{a} \leftarrow \vec{\bm{a}}[0] \in \Field$; $V_G \leftarrow \vec{\bm{G}}[0] \in \Group_1$; \;
	\label{line:reduceHata}
		$\Verifier$ computes:\\
	\Indp
		$\vec{\bm{\xi}}_G \leftarrow (1, x_{(\log_2 n +1)}) \kronecker \cdots \kronecker (1, x_{(2\log_2 n)}) \in \Field^n$\;
		\Indm
	$\Verifier$ checks:
		$C\, \gadd \, \gsum_{j=1}^{2\log_2(n)} (x_{(j)} L_{(j)} \gadd x_{(j)}^{-1}R_{(j)}) \EqOrNot$ 
		$ 
        [\langle \vec{\bm{\xi}}_G ,\,\vec{\bm{l}}\rangle \cdot
        \langle \vec{\bm{\xi}}_H ,\, \vec{\bm{r}} \rangle
        \cdot (\hat{a} G_0 H_0)]
		\gadd
		(\hat{a} V_G V_H)
		$\;
		\label{line:verifierPCS}
	$\Prover$ provides proofs (see Appendix~\ref{appx:g2acc}) to let $\Verifier$ verify:
		$V_G \EqOrNot \langle \vec{\bm{\xi}}_G , \vec{\bm{G}}\rangle$; $ V_H \EqOrNot \langle \vec{\bm{\xi}}_H, \vec{\bm{H}} \rangle$\;
        \label{line:gtbulletfinalcheck}
\end{algorithm}

\mysmallskip{\bf Square-root group operations in $\Group_T$-Bulletproofs.}
We observe that 
$\Group_T$-Bulletproofs iteratively fold the $n^2$-vector $\vec{\bm{G}} \kronecker \vec{\bm{H}}$.
To achieve square-root 
%pairings and 
group operations,
%in Bulletproofs, 
we propose replacing these $\Group_T$ foldings with more efficient $\Group_1$ or $\Group_2$ foldings on shorter $n$-vectors,
exploiting the following tensor product property:
\begin{align*}
& (\vec{\bm{G}}
	\kronecker 
	\vec{\bm{H}})_\mathsf{L} \,
\gadd \, 
x_{(1)} 
	(\vec{\bm{G}}
	\kronecker
	\vec{\bm{H}})_\mathsf{R} =
	(\vec{\bm{G}}_\mathsf{L} \, \gadd \, x_{(1)}\vec{\bm{G}}_\mathsf{R}) \kronecker \vec{\bm{H}}.
\end{align*}
While the LHS folding has $O(n^2)$ complexity, 
the RHS folding requires only $O(n)$.
By utilizing this optimization, $\Group_T$-Bulletproofs achieve the same asymptotic efficiency as Dory.
The concrete numbers of pairings and group operations are similar to BMMTV21.

This $\Group_T$-Bulletproofs are presented in Algorithm~\ref{alg:pcs}.

\mysmallskip
\noindent
{\bf Trusted setup only for the final check.}
Algorithm~\ref{alg:pcs} is compatible with both transparent and trusted setups, depending on how
$\vec{\bm{G}}$ and $\vec{\bm{H}}$ are generated.
A trusted setup is required only for the final public multi-exponentiation checks
$V_G = \langle \vec{\bm{\xi}}_G , \vec{\bm{G}} \rangle \in \Group_1$ and
$V_H = \langle \vec{\bm{\xi}}_H , \vec{\bm{H}} \rangle \in \Group_2$
at Line~\ref{line:gtbulletfinalcheck}.
Under a transparent setup, these equations can be verified in logarithmic time
(see~\cite[\S 4.5]{iacr/cong25}),
whereas a trusted setup enables constant-time verification, as discussed in
Appendix~\ref{appx:g2acc}.

\mysmallskip
\noindent
{\bf Bucket-wise opening.}
For the expanded vector produced by a BVA, 
Algorithm~\ref{alg:pcs} can be executed without materializing the BVA.
In Algorithm~\ref{alg:pcs}, the only computations involving the witness $\bm{a}$
are the matrix projections in Lines~\ref{line:computeal} and~\ref{line:computear}.
These matrix projections can be performed bucket by bucket
without materializing padded coordinates.
Conceptually, the expanded vector is reshaped into a square matrix.
The prover iterates only over the real buckets, computes the position of each element within the expanded vector, multiplies the element by the corresponding entry in $\vec{\bm{l}}$ or $\vec{\bm{\xi}_H^{-1}}$, respectively, and accumulates the result to obtain $\vec{\bm{l}}^\top \bm{a}$ and $\bm{a}\vec{\bm{\xi}_H^{-1}}$.
All other computations in Algorithm~\ref{alg:pcs} are independent of $\bm{a}$.


\subsection{Verifier acceleration via trusted setup}
\label{appx:g2acc}
In Line~\ref{line:gtbulletfinalcheck} of Algorithm~\ref{alg:pcs}, the verification requires computing two multi-exponentiations
$V_G = \langle \vec{\bm{\xi}}_G , \vec{\bm{G}} \rangle \in \Group_1$ and
$V_H = \langle \vec{\bm{\xi}}_H, \vec{\bm{H}} \rangle \in \Group_2$,
with public tensor-structured vectors $\vec{\bm{\xi}}_G$ and $\vec{\bm{\xi}}_H$.
zkMatrix~\cite[\S 3]{asiaccs/CongYY24} suggests that using the trusted setup,
the prover can provide additional proof for the $\Group_1$ multi-exponentiation $V_G = \langle \vec{\bm{\xi}}_G, \vec{\bm{G}} \rangle \in \Group_1$ via two extra $\Group_1$ multi-exponentiations.



To verify the $\Group_2$ multi-exponentiation $V_H = \langle \vec{\bm{\xi}}_H, \vec{\bm{H}} \rangle \in \Group_2$.
Define a polynomial:
\[
	\phi(s): = \xi_{H1} + \xi_{H2}s^q + \cdots + \xi_{Hn} s^{q(n - 1)}.
\]
For a random challenge $s \ChooseRandom \Field^*$, compute the coefficients of $\hat{s}^q$ in:
\[
	\frac{\phi(s) - \phi(\hat{s})}{s^q - \hat{s}^q} = \zeta_1 + \zeta_2 \hat{s}^q + \cdots + \zeta_{n - 1} \hat{s}^{q(n 
 - 2)},
\]
collecting these coefficients into $\vec{\bm{\zeta}}_H \in \Field^{n - 1} $.

Using the trusted setup, the prover computes:
\[
	V_H' \leftarrow \vec{\bm{\xi}}_H \bullet \vec{\bm{H}}'' \in \Group_2, \quad
	W_H \leftarrow \vec{\bm{\zeta}}_H \bullet \vec{\bm{G}}' \in \Group_1.
\]
The verifier checks two pairing equations:
\[
	e(\hat{G}, V_H') = e(G'', V_H), \quad
	e(\hat{G}, \phi(s)\hat{H} \gminus V_H) = e(W_H, s^q \hat{H} \gminus H').
\]
This check runs in constant time and is sufficient to ensure that $V_H = \langle \vec{\bm{\xi}}_H, \vec{\bm{H}} \rangle$.

Using this technique, the overall proof size and verifier time of
Algorithm~\ref{alg:pcs} are logarithmic.

Knowledge soundness of this approach is demonstrated similarly to the $\Group_1$ multi-exponentiation method in ~\cite[Appendix A.1]{asiaccs/CongYY24} and is therefore omitted.
